% Mathematical note for the root agent; intended for inclusion in the report.
% Baseline: ProveIt 443cff49563457a48ee0f663012d56c4bdb58fa0.
% Requires amsmath, amssymb; no custom theorem environments are required.

\section{Factoring every successor iteration into one unit step}

The new restriction is that the arithmetic evaluation rule needs only
iterations whose counter is the fixed numeral $1=S0$.  It has no variable
counter and does not reconstruct or compare a predecessor iterator.
The price is an explicitly specified, purely syntactic factorization in
the term representation.  No iterator body is executed by this
factorization.  This section proves that the change preserves the entire
consequence relation, with arbitrary retained assumptions.

\subsection{Raw and canonical terms}

Raw terms are
\[
 t ::= x\mid 0\mid St\mid I(n,a;u.h),
\]
where $u$ is bound only in $h$.  Let
\[
 U(a;u.h):=I(1,a;u.h)
\]
be an abbreviation, not a new term constructor.  Copies of $h$ in
different iterator nodes have independent binders; bound variables may
always be renamed before substitution.

Define a smart constructor $K$ on already canonical arguments by
\begin{align*}
 K(0,a;u.h)&=a,\\
 K(Sn,a;u.h)&=U(K(n,a;u.h);u.h),\\
 K(n,a;u.h)&=I(n,a;u.h)
       &&\text{if the root of $n$ is neither $0$ nor $S$.}
\end{align*}
The second clause is an algorithmic call on the proper subterm $n$ of
$Sn$.  It does not recursively normalize its output.  In particular,
$K(1,a;u.h)=U(a;u.h)$, and this unit application is retained.

Define $N$ by structural recursion:
\begin{align*}
 N(x)&=x,&N(0)&=0,&N(St)&=S N(t),\\
 N(I(n,a;u.h))&=K(N(n),N(a);u.N(h)).
\end{align*}
Its fixed points are precisely the recursively well-formed terms in
which every iterator counter is either literally $1$, or has a root
other than $0$ and $S$.  Thus no zero-counter iterator remains, and the
only successor-counter iterator is a unit application.

For example,
\[
 N(I(SSx,a;u.h))
 =U\bigl(U(I(x,N(a);u.N(h));u.N(h));u.N(h)\bigr),
\]
when $x$ is a variable.  The displayed nested uses of $u$ bind their own
bodies only.  Neither $h[a/u]$ nor any other execution of $h$ appears in
the normalization algorithm.

Extend $N$ to formulas by changing only numerical terms, and to clauses
and contexts componentwise.  A substitution in the canonical calculus
means capture-free raw substitution followed by $N$.

\subsection{Termination and confluence of the corresponding rewriting}

One can also describe the normal forms by the following two contextual
rewrites, allowed inside every term and under every binder:
\begin{align*}
 I(0,a;u.h)&\longrightarrow a, \tag{Z}\\
 I(Sn,a;u.h)&\longrightarrow U(I(n,a;u.h);u.h)
       &&(n\text{ is not literally }0). \tag{F}
\end{align*}
The side condition on (F) is essential for this rewrite description:
without it, rewriting $U(a;u.h)$ would expand it to
$U(I(0,a;u.h);u.h)$, which (Z) contracts back.  The smart-constructor
algorithm above is already terminating and needs no such side-condition
test in its successor clause.

Here is a direct termination proof, including rewriting in bodies and
the duplication of those bodies.  Give terms positive integer weights
$w$ by
\[
 w(x)=w(0)=1,\qquad w(St)=w(t)+1,
\]
and, writing $A=w(a)$ and $H=w(h)$, set
\[
 w(I(n,a;u.h))=
 \begin{cases}
 A+H+1,&n\text{ is literally }1,\\
 (w(n)+1)(A+H+1),&\text{otherwise}.
 \end{cases}
\]
Rule (Z) decreases the weight from $2(A+H+1)$ to $A$.
For (F), its left-hand counter is not literally $1$, so its weight is
$(w(n)+2)(A+H+1)$.  If $n$ is not literally $1$, the right-hand side
has weight
\[
 (w(n)+1)(A+H+1)+H+1,
\]
and the decrease is exactly $A>0$.  If $n$ is literally $1$, the
right-hand side has weight $A+2H+2$, whereas the left-hand side has
weight $4(A+H+1)$, again strictly greater.

These decreases are preserved in every surrounding context.  The
weight is strictly increasing in initializer and body weights.  In the
counter position the only exceptional case is a rewrite making the
counter literally $1$; it changes the multiplicative clause to the
smaller additive clause.  A counter already equal to the irreducible
literal $1$ cannot rewrite away from it.  Successor contexts are
strictly monotone.  Induction on contexts therefore shows that every
rewrite strictly decreases a positive integer.  This proves strong
normalization.

Both rewrites preserve $N$.  For (F), use the identity
\[
 K(Sn,a;u.h)=U(K(n,a;u.h);u.h),
\]
which is true even when $n=0$.  It gives
\begin{align*}
 N(I(Sn,a;u.h))
 &=U(K(N(n),N(a);u.N(h));u.N(h))\\
 &=N(U(I(n,a;u.h);u.h)).
\end{align*}
Rule (Z) is immediate.  A structural induction extends this invariance
to contextual rewrites.  Every irreducible term is canonical, and hence
is fixed by $N$; conversely, every canonical term is irreducible.
Consequently every term has exactly the normal form $N(t)$, and any
two reduction sequences can be joined at this normal form.  The system
is confluent.

All the constructions are effective on finite syntax.  With an ordinary
primitive-recursive coding of trees, $w$, individual rewrites, and a
deterministic choice of the next redex are primitive recursive.  At most
$w(t)-1$ rewrites are possible.  Thus normalization has a primitive
recursive implementation, independently of the smart-constructor
implementation.  This argument does not assume a uniform primitive
recursive evaluator for numerical programs: no body is substituted into
an application during normalization.

\subsection{Substitution coherence}

For every capture-free substitution $\sigma$,
\[
 N(t[\sigma])=N\bigl(N(t)[N\circ\sigma]\bigr). \tag{SC}
\]
Here substitution is lifted in the usual way under the binder $u$.
One proof separates two elementary facts.  First, replacing every
substitution image by its $N$-normal form does not change the normal
form of the substituted term: apply the preceding contextual rewrite
invariance to each inserted image, with fresh bound names.
Second, if $t\longrightarrow t'$ is one rewrite, then
$N(t[\sigma])=N(t'[\sigma])$.  The zero rule is preserved literally.
For the factorization rule the substitution can turn its nonzero
syntactic parameter $n$ into a term normalizing to zero; nevertheless
the displayed identity for $K(Sn,a;u.h)$ is valid also at zero, so the
normal forms are still identical.  Follow a finite reduction of $t$ to
$N(t)$ and combine these two facts.  This proves (SC) without assuming
that the conditional rewrite relation itself is closed under arbitrary
substitution.

It follows that canonical substitution is associative:
\[
 N\bigl(N(t[\sigma])[\tau]\bigr)
 =N\bigl(t[x\mapsto N(\sigma(x)[\tau])]\bigr).
\]
Normalization is idempotent, commutes with renaming, and introduces no
new free variables.  Formula and clause versions follow by structural
induction, with the appropriate lifted substitutions.

\subsection{The sole arithmetic inference pattern}

Use the following two schemes on canonical terms:
\[
 \overline{\Gamma\vdash t=t}\quad\mathrm{Refl},
 \qquad
 \frac{\Gamma\vdash a=b\qquad
       \Gamma\vdash U(a;u.h)=r}
      {\Gamma\vdash r=N(h[b/u])}\quad\mathrm{Unit}.
\]
Unit checks the fixed counter $1$ and compares its initializer directly
with the left side of its first premise.  It has no arbitrary counter,
no construction of a predecessor iterator, and no supplied equality
template separate from the displayed application body.

For $J_s:=U(s;u.u)$, the same explicit pivot as in the baseline is now
made entirely of unit applications.  From $s=t$ and $s=r$ obtain
\[
 \frac{s=t\qquad
       \displaystyle\frac{s=r\qquad J_s=J_s}{J_s=r}\ \mathrm{Unit}}
      {r=t}\ \mathrm{Unit}. \tag{Pivot}
\]
Symmetry and transitivity follow by reflexivity and Pivot.  For a fixed
template $h$, put $H_s=U(s;u.h)$.  Applying Unit once with $s=s$ and
once with $s=t$, using $H_s=H_s$ as the other premise, gives
$H_s=N(h[s/u])$ and $H_s=N(h[t/u])$.  Pivot gives term congruence.
One further Pivot recovers crossed one-template transport
\[
 s=t,\quad N(h[s/u])=r\quad\vdash\quad r=N(h[t/u]).
\]
No induction or numerical case distinction is used in this equality
argument.

The normalized generic successor source is definitionally
\[
 K(Sc,a;u.h)=U(K(c,a;u.h);u.h).
\]
Therefore any baseline Step instance
\[
 I(c,a;u.h)=t,\quad I(Sc,a;u.h)=r
       \quad\vdash\quad r=h[t/u]
\]
is translated into exactly one Unit instance.  Its first premise
supplies the equality of the unit application's initializer with $t$.
In particular, reflexivity as both premises derives normalized
successor computation for every counter $c$.

\subsection{The complete five-scheme presentation}

Keep the baseline Eq/$D$ formulas and reusable clauses:
\[
 A::=(s=t)\mid D_{x,y}B(x,y),\qquad
 [A;y.B(y)]^*:=A^*\to\forall y\,B(y)^*,
\]
with
\[
 (D_{x,y}B)^*:=\exists x\bigl(B(x,0)^*\to\forall y\,B(x,y)^*\bigr).
\]
The bound $y$ of a clause does not bind its antecedent.  Write
$\operatorname{Assm}(A)=[0=0;y.A]$ with $y$ fresh for $A$.
Besides Refl and Unit, retain exactly these three schemes:
\begin{align*}
 \mathrm{Use}:\quad&
 \frac{[A;y.B]\in\Gamma\qquad\Gamma\vdash A}
      {\Gamma\vdash N(B[t/y])},\\[2mm]
 \mathrm{DE}:\quad&
 \frac{\Gamma\vdash D_{x,y}B\qquad
       \Gamma,[N(B[b/x,0/y]);y.N(B[b/x])]\vdash A}
      {\Gamma\vdash A},\\[2mm]
 \mathrm{DI}:\quad&
 \frac{\Gamma,\operatorname{Assm}(N(B[t/x,\eta/y]))
           \vdash N(B[t/x,Sz/y])}
      {\Gamma\vdash D_{x,y}B}.
\end{align*}
In DE, $b$ is fresh for the outer context and conclusion.  In DI, $z$
is fresh for the outer context, witness $t$, and passive parameters,
and $\eta=z$ exactly when $B$ is a literal equation and $t$ is
literally zero; otherwise $\eta=0$.  All the terms and formulas are
well scoped before opening a branch.  Outer contexts are weakened
when a fresh variable is introduced.  These are the baseline logical
rules, with the new canonical substitution; no extra logical or
arithmetic axiom is implicit.

\subsection{Why the introduction mode survives the new normalization}

Let $N_0$ be zero-only normalization.  For every raw term $t$,
\[
                 N(t)=0\quad\Longleftrightarrow\quad N_0(t)=0.
                 \tag{Zero}
\]
Prove this by structural induction.  A variable or successor cannot
normalize to zero.  For an iterator, either its normalized counter is
zero, in which case the result is its normalized initializer, or its
normal form remains iterator-headed.  Apply the induction hypotheses
to the counter and initializer.  The new factorization never removes a
unit application or exposes its body.

In particular, if $t$ is already zero-normal, $N(t)$ is literally zero
exactly when $t$ is literally zero.  Formula root constructors are also
preserved.  Hence the witness-directed DI mode is preserved when the
existing zero-normal Step calculus is translated by $N$.  The new
factorization creates no new witness-mode problem.  The older raw-to-
zero-normal transformation has its already documented mode repair,
which can be composed with this transformation if raw terms are the
starting presentation.

\subsection{Exact equivalence and the numerical interpretations}

Let $C_0$ denote the baseline zero-normal five-scheme clause calculus
with Refl and Step, and let $C_1$ be the present calculus.  Translating
each term and formula by $N$ gives a rule-by-rule simulation
\[
       C_0:\Gamma\vdash A\quad\Longrightarrow\quad
       C_1:N\Gamma\vdash NA.
\]
Refl is immediate.  Step is the single Unit instance just exhibited.
Use and DE follow from substitution coherence.  DI follows from that
coherence and (Zero), which preserves its mode and freshness.  This is
a translation of general derivations, not just closed arithmetic
calculations.

For the reverse simulation, interpret the canonical syntax in the
baseline ordinary scalar theory.  Both (Z) and (F) express provable
equations.  For (F), the two sides equal $h[I(n,a;u.h)/u]$ by successor
computation, with the right side using the unit and zero computations.
Thus every term has the same provable value as its $N$-normal form.
When normalization changes a body under its iterator binder, one uses
the baseline's derived body-congruence lemma: pointwise equality of
two bodies yields equality of their iterates by literal-equation
induction on the counter.  This dependency must be counted; ordinary
term congruence alone does not identify arbitrary bodies under a
binding operator.

Unit is the $c=0$ instance of baseline Step, followed when necessary
by the provable equality between a raw substitution result and its
normal form.  Its ordinary reading is also immediate from one-step
computation and equality.  The three logical rules have their baseline
readings; replacing raw substitution instances by normal forms changes
them only by provable equivalences.  Consequently the reverse
simulation holds as well.  Equivalently, for every canonical sequent,
the two presentations have precisely the same consequences; for a raw
sequent, first apply $N_0$ and then $N$.

The baseline's whole-domain bi-interpretation with full first-order
iPRA therefore transfers unchanged.  Normal forms concern expression
syntax only.  They do not quotient the numerical domain, restrict it
to denoted elements, or replace actual equality by a numerical test.
All numerical operations retain their baseline primitive-recursive
meanings.  Both numerical interpretation domains are the whole domain,
and both composite isomorphisms have graph $x=y$.  On formulas,
normalization is provably equivalent to the original formula by the
value equalities and the ordinary logical congruence rules.  The
designated contradiction $0=S0$ is unchanged.  Thus bi-interpretability
and equiconsistency are both preserved.

\subsection{Scope of the simplicity claim}

This is a stricter arithmetic inference interface: only the unit
instance of Step remains, and its premise matching uses the
initializer directly.  The number of inference schemes is still five.
The factorizing constructor and its scope-aware implementation are
part of the complete specification and must be included in every
source-size comparison.  A reduction in the inference rule by itself
does not establish a reduction in total checker size.  No absolute
minimality, priority beyond the inspected baseline, or proof-assistant
verification is asserted here.
